\chapter{Implementation}

\section{Typed domain boundary}

The implementation is organized as small standard-library modules. \texttt{EOScene} is a frozen
dataclass containing an identifier, platform, acquisition time, cloud cover, GeoJSON footprint,
optional bounding box, EPSG code, and an asset map. The provider adapter parses an untrusted STAC
Item into this type only after checking required fields, coordinate structure, cloud-cover range,
asset URLs, and temporal syntax. A frozen value object prevents downstream catalog and report code
from accidentally mutating the record it is measuring.

The transformation boundary is deliberately one-way:
\[
\text{provider JSON}\xrightarrow{\text{validate}}\text{STAC item}
\xrightarrow{\text{normalize}}\text{EOScene}
\xrightarrow{\text{persist}}\text{catalog and provenance}.
\]
Malformed input is a permanent input error. It is not retried, because repeating an invalid payload
cannot improve its schema.

\section{Bounded asynchronous ingestion}

\texttt{ConcurrentEOIngestor} combines asynchronous provider work with synchronous persistence.
The caller supplies an iterable of items, a transformation function, a provenance sink, and a
maximum worker count \(k\). A semaphore with capacity \(k\) is acquired immediately before
provider/transform work and released in a \texttt{finally} block. Thus, for active tasks
\(A(t)\),
\[
0\leq A(t)\leq k\qquad\forall t.
\]
This invariant prevents a large provider response from creating an unbounded number of network
tasks or file operations.

The scheduling pattern is summarized below.

\begin{verbatim}
semaphore = asyncio.Semaphore(max_workers)

async def ingest_one(item):
    async with semaphore:
        scene = transform(item)
        await provenance.write(scene)
        await asyncio.to_thread(catalog.upsert, scene)
        return scene

results = await asyncio.gather(
    *(ingest_one(item) for item in items),
    return_exceptions=False,
)
\end{verbatim}

The real implementation records counts, elapsed time, payload bytes, retries, and failures in
\texttt{IngestionMetrics}. \texttt{asyncio.to\_thread} keeps blocking filesystem and SQLite calls
off the event loop, while an asynchronous lock serializes append operations to the shared
provenance manifest. A worker-limit stress test observes the peak active count and verifies that
all valid scenes are eventually persisted.

\section{Provenance and content addressing}

Before hashing, a payload is serialized as canonical UTF-8 JSON with deterministic key ordering and
compact separators. For byte sequence \(x\), the content address is
\[
h=\operatorname{SHA256}(x),
\]
represented as a lowercase hexadecimal digest. The raw payload is stored below a digest-derived
path, while an append-only JSON Lines record stores the digest, source URI, provider identifier,
timestamp, and retry metadata. Since identical canonical payloads have identical \(h\), repeated
observations can be detected without trusting a mutable filename.

The manifest is an audit log rather than a mutable status table. A successful write therefore has
two observable products: the addressed raw document and a manifest line describing how it entered
the system. The sink writes through a temporary file and atomic replacement for the raw payload,
then appends a complete line under its lock. This ordering avoids exposing a manifest reference to
a missing raw object. Provenance cleanup inspects references before deleting stale files and never
targets the frozen experiment report.

\section{Quality profiling and catalog persistence}

The metadata-only quality profiler evaluates geometry validity, required asset completeness, and
cloud risk without downloading raster data. A scene is usable when required structural predicates
hold; the resulting profile records missing assets, risk flags, and a bounded quality score.
Catalog SQL applies indexed platform, temporal, cloud, and bounding-box filters first. Exact
polygon checks then remove false positives from the rectangular pre-filter.

\texttt{SpatialCatalog} uses a primary key on scene identity and an idempotent upsert. A write is
committed in a transaction, and readers use short-lived queries under a reentrant connection lock.
WAL mode permits readers to continue while a writer appends changes. The mixed-reader/writer
regression test accepts any committed snapshot between the initial and final counts, checks that
no reader observes invalid state, and reopens the database to verify durability.

\section{NDVI analytical slice}

The raster module intentionally stops at range retrieval and a small numerical transformation.
For a pixel with near-infrared reflectance \(N\) and red reflectance \(R\), the Normalized
Difference Vegetation Index is
\[
\operatorname{NDVI}(N,R)=\frac{N-R}{N+R}.
\]
The denominator is tested before division:
\[
\operatorname{NDVI}(N,R)=
\begin{cases}
0, & N+R=0,\\[2mm]
\dfrac{N-R}{N+R}, & \text{otherwise}.
\end{cases}
\]
The zero result is a deliberate finite sentinel for undefined zero-signal pixels; it avoids a
floating-point exception and keeps array-free standard-library processing deterministic. The
\texttt{RasterAssetFetcher} requests byte ranges with an HTTP \texttt{Range} header and leaves
full raster decoding to a future dependency-backed layer.

\section{Multiscale downsampling and grid alignment}

The multiscale raster extension operates on a rectangular numerical band represented as nested
Python tuples. For reduction factor \(k>0\), output cell \((r,c)\) contains the source block
\[
\mathcal{B}_{r,c}=
\{x_{y,z}\mid rk\leq y<\min((r+1)k,H),\
ck\leq z<\min((c+1)k,W)\}.
\]
The brightness-preserving reduction is the arithmetic mean
\[
D_{r,c}=\frac{1}{|\mathcal{B}_{r,c}|}
\sum_{x\in\mathcal{B}_{r,c}}x.
\]
The denominator is the actual block cardinality, including at image edges. Therefore a constant
input field \(x_{y,z}=v\) remains \(D_{r,c}=v\) for every complete and partial block. The
implementation never divides an accumulated sum by \(k^2\) after already averaging or normalizing
the samples, avoiding artificial darkening and preserving the input dynamic range.

\texttt{downsample_scene_bands} applies the same factor to named spectral grids such as B04 and
B08, preserving their spatial correspondence and relative calibration. The benchmark helper
records input/output pixel counts, monotonic elapsed time, pixels per second, and Python peak
allocation measured with \texttt{tracemalloc}. It is a reproducible memory/throughput indicator,
not a measurement of native raster-library memory or disk I/O.

\texttt{AffineGridTransform} supplies the corresponding numerical alignment contract. Given
EPSG:4326 bounding box \((\lambda_{\min},\phi_{\min},\lambda_{\max},\phi_{\max})\) and dimensions
\((W,H)\), its forward mapping is
\[
(\lambda,\phi)=
\left(\lambda_{\min}+\frac{c}{W}(\lambda_{\max}-\lambda_{\min}),
\phi_{\min}+\frac{r}{H}(\phi_{\max}-\phi_{\min})\right),
\]
with the algebraic inverse mapping coordinates back to fractional pixel positions. This utility
supports deterministic local-grid alignment for synthetic scenes; it does not claim datum-aware
UTM reprojection and deliberately leaves ellipsoidal projection and GeoTIFF decoding to future
optional processing layers.

\section{OAuth2 recovery workflow}

\texttt{CDSETokenProvider} obtains client-credentials tokens through an injectable transport.
Environment variables provide the token URL, client ID, client secret, and STAC URL. The provider
stores the token acquisition instant and refreshes with a skew \(\delta\) before expiry. With
provider lifetime \(L\), a cached token is valid iff
\[
t_{\mathrm{now}} < t_{\mathrm{issued}}+ \max(0,L-\delta).
\]
The authorization header is constructed only at the request boundary, so the secret and token are
not written to reports or logs.

Transient failures use bounded exponential backoff. For retry attempt \(n\), base delay \(d\), cap
\(D\), and optional provider retry delay \(r\), the sleep duration is
\[
\tau_n=\min(D,d\,2^n),
\qquad
\tau_n'=\max(\tau_n,r)
\]
when \(r\) is supplied. The loop retries timeouts and HTTP 429 responses only; an invalid JSON
shape, missing access token, or non-positive lifetime raises a protocol error immediately.

\begin{verbatim}
for attempt in range(policy.max_attempts):
    try:
        payload = await transport.post_form(token_url, form)
        token = validate_token_payload(payload)
        cache.store(token, clock())
        return token
    except (TimeoutError, RateLimited) as error:
        if attempt + 1 == policy.max_attempts:
            raise
        await asyncio.sleep(policy.delay(attempt, error.retry_after))
\end{verbatim}

An expired cached token is discarded before the next request. If a provider returns an
authentication failure after caching, the client invalidates the token and performs one bounded
refresh path. Tests cover expiry, 429 responses, timeouts, exhausted retries, and malformed token
payloads so that recovery behavior is observable rather than inferred from a successful path.

\section{Reports and live experiment operations}

The reporting layer serializes metrics, catalog counts, quality distributions, spatial query
latency, recovery ratio, and payload footprint to sorted JSON and human-readable Markdown. The
offline benchmark uses deterministic fixtures and remains frozen at 256 scenes. The live runner
resolves CDSE configuration only when all required variables are present; otherwise it uses the
fixture transport. Each live run writes a date-stamped JSON/Markdown pair and refuses same-day
overwrite. A dry-run-by-default maintenance command prunes only dated live pairs older than the
retention threshold.

This separation gives the implementation two operational modes: a reproducible, credential-free
mode for CI and teaching, and an authenticated mode for empirical provider measurements. Both
produce the same report schema, allowing live latency and response distributions to be compared
with the frozen baseline without mutating it.
